% 本文件由 LaTeX 排版 Agent 根据有效 translation.json 自动生成。
% author=Councils; work_id=3801; title=尼西亚第一届大公会议（公元325年）

\chapter{尼西亚第一届大公会议（公元325年）}

% structure: p0001 source=h1 target=omitted reason=page-title-already-rendered-by-parent

% structure: p0002 source=h2 target=section reason=relative-heading-rank; base=h2

\section{尼西亚信经}

尼西亚会议公佈了本信经。\par

尼西亚会议之信条声明。\par

我们信唯一的天主，全能的圣父，有形与无形万物的创造者；并信唯一的主耶稣基督，天主子，由圣父所生，父的独生子，出自父的本体，天主出自天主，光明出自光明，真天主出自真天主，受生（\OriginalTerm{受生}{\GreekText{γεννηθέντα}}）而非受造，与圣父同一本体（\OriginalTerm{同质}{\GreekText{ὁμοούσιον}, consubstantialem}）。万物藉祂而造，无论在天上或在地上。祂为了我们人类并为了我们的救恩，从天降下，成了血肉，成为人。祂受难，第三日复活，升了天。祂还要降来审判生者死者。我们也信圣神。凡说\Quote{曾有一时，圣子不存在}（\OriginalTerm{曾有一时，圣子不存在}{\GreekText{ἤν} \GreekText{ποτε} \GreekText{ὅτε} \GreekText{οὐκ} \GreekText{ἦν}}），或说祂在受生之前不存在，或说祂由无中造成，或说祂与圣父本体或本质不同，或说祂是受造物，或说祂可变化或可转变的——凡如此说的，公教与宗徒的教会都予以绝罚。\par

% structure: p0006 source=h2 target=section reason=relative-heading-rank; base=h2

\section{教规}

% structure: p0007 source=h3 target=subsection reason=relative-heading-rank; base=h2

\subsection{教规一}

若有人因病由医者施行手术，或被蛮人阉割，可仍留于圣职人员中；但若在健康时自阉，则当如此者若已列于圣职中，应停止其职务，且此后不得升任任何圣职。然而，此规定显然指向那些故意行事并擅自阉割者；至于被蛮人或其主人造成阉人者，且在其他方面堪当，则本教规准其加入圣职。\par

% structure: p0009 source=h3 target=subsection reason=relative-heading-rank; base=h2

\subsection{教规二}

\Emph{既然}因必要或个人的催促，许多事违反教会教规而行，以致刚由异教归信者，受训未久，即被引至神性洗礼池，且一经领洗，即擢升为主教或司铎；我们认为此后不宜再做此类事。因为连望教者本身也需要时间，领洗后亦需更长期的考验。因圣经明言：\Quote{\BibleQuote{不可拣选新奉教的人，免得他骄傲自大，而陷于魔鬼所受的审判和罗网。}}（\BibleRef{弟前 3:6}）但若日后此人被发现有任何感官之罪，且经两三个证人证实，则当令其停止圣职。凡违反这些规条者，将危及自己的圣职地位，因其擅敢违背此大公会议。\par

% structure: p0011 source=h3 target=subsection reason=relative-heading-rank; base=h2

\subsection{教规三}

大公会议严厉禁止任何主教、司铎、执事或任何圣职人员，与\Emph{\OriginalTerm{同居女子}{subintroducta}}同住，除非是母亲、姊妹、姑姨或任何无可疑之人。\par

% structure: p0013 source=h3 target=subsection reason=relative-heading-rank; base=h2

\subsection{教规四}

主教应由本省所有主教任命；但若因紧急需要或距离遥远而难以做到，则至少应有三位主教聚集，并征得缺席主教的同意且以书面传达后，方可举行祝圣。但在每一省，所行之事应由都主教批准。\par

% structure: p0015 source=h3 target=subsection reason=relative-heading-rank; base=h2

\subsection{教规五}

关于在各省被绝罚者，无论圣职人员或平信徒，应遵循教规：被某些人驱逐者，不得被他人重新接纳。但仍应查明，是否因主教的挑剔、争辩或类似不友善态度而被绝罚。为使此事得到适当调查，决定每省每年应举行两次教务会议，俾使本省全体主教聚集，详细审查此类问题，使那些确实得罪其主教者，可被众人视为因正当理由被绝罚，直至主教全体会议认为适宜对他们宣布较轻判决。这些教务会议一次应在四旬期之前举行（使一切怨气消除后，可将纯洁的祭品奉献给天主），第二次应在秋季举行。\par

% structure: p0017 source=h3 target=subsection reason=relative-heading-rank; base=h2

\subsection{教规六}

愿埃及、利比亚和彭塔波利的古老习俗得以维持，即亚历山大里亚主教对所有这些地区有管辖权，因罗马主教亦照此惯例。同样，在安提约基雅及其他省份，各教会应保留其特权。此规定普世适用：若有人未经都主教同意而成为主教，大公会议宣布其人不应为主教。但若两三位主教因自然好争论而反对其余主教的共同投票，而多数投票合理且符合教会法律，则多数之选择应生效。\par

% structure: p0019 source=h3 target=subsection reason=relative-heading-rank; base=h2

\subsection{教规七}

鉴于习俗与古老传统维护了埃利亚（\Emph{即}耶路撒冷）主教应受尊敬，故应保存其大都会的应有尊严，而埃利亚主教享有次位的尊荣。\par

% structure: p0021 source=h3 target=subsection reason=relative-heading-rank; base=h2

\subsection{教规八}

关于那些自称为\OriginalTerm{洁净派}{Cathari}的人，如果他们归向天主教会和宗徒的教会，伟大而神圣的主教会议规定，凡已受圣秩者，应继续留在圣职人员中。但首要的是，他们必须以书面宣认遵守并遵行天主教会和宗徒的教会的信理；尤其是他们将与再婚者以及与那些在迫害中失足并被规定补赎期限和复和时限的人共融，从而在一切事上遵行天主教会信理。因此，无论在村庄或城市，若发现所有受圣秩者仅属此类，则让他们留在圣职人员中，并保持其原有品级。但如果他们归向已有天主教会的主教或司铎的地方，则显然该教会的主教应拥有主教之尊；而由那些被称为洁净派的人所任命的主教，应保有司铎的品级，除非主教认为适宜给他分享主教衔的荣誉。或者，若此安排不满意，则主教应为他提供一个\OriginalTerm{乡村主教}{Chorepiscopus}或司铎的职位，以便他明显可见地属于圣职人员，且在同一城市中不致有两位主教。\par

% structure: p0023 source=h3 target=subsection reason=relative-heading-rank; base=h2

\subsection{第九条 教规}

若有任何司铎未经审查即被擢升，或经审查后他们承认了罪行，而有人违反教规（不顾其认罪）仍为他们覆手，则教规不予接纳；因为天主教会只要求无可指摘者。\par

% structure: p0025 source=h3 target=subsection reason=relative-heading-rank; base=h2

\subsection{第十条 教规}

若有任何失足者因授秩者的无知或甚至明知而被授予圣秩，这不影响教会的教规；因为他们一经发现，即应被撤职。\par

% structure: p0027 source=h3 target=subsection reason=relative-heading-rank; base=h2

\subsection{第十一条 教规}

关于那些非因强迫、非因财产被剥夺、非因危险或类似情况而跌倒者（例如在里基尼乌斯暴政期间所发生的），主教会议宣布：虽然他们不配得到宽恕，但仍应以仁慈对待。凡曾是信友者，若诚心悔改，应在听道者中度过三年；作俯伏者七年；并与百姓在祈祷中共融两年，但不奉献祭品。\par

% structure: p0029 source=h3 target=subsection reason=relative-heading-rank; base=h2

\subsection{第十二条 教规}

凡蒙受恩宠而表现起初热忱、抛弃了军衔腰带，但后来又像狗一样回到自己的呕吐物者（以致有些人花费金钱，藉礼物重新获得军职）：这些人应在听道者中度过三年，然后作俯伏者十年。但在所有这些情形中，必须仔细考察他们的意向以及他们的悔改是怎样的。因为凡以行为而非假装，以恐惧、眼泪、恒心与善工证明自己悔改的人，当他们完成了听道者的指定时间后，可以适当地在祈祷中共融；此后，主教可以对他们作出更仁慈的判定。但那对此漫不经心、认为不进教堂的形式就足以使他们悔改的人，必须完成整个期限。\par

% structure: p0031 source=h3 target=subsection reason=relative-heading-rank; base=h2

\subsection{第十三条 教规}

关于临终者，古老的教规法律仍应维持，即：若有人处于死亡边缘，不可剥夺他最后且最必要的\OriginalTerm{临终圣体}{Viaticum}。但如果有人领受了圣体后生命获救恢复健康，则让他仅与那些在祈祷中共融的人在一起。但总的来说，在任何垂死之人请求领受圣体圣事的情况下，主教经审查后应给予他。\par

% structure: p0033 source=h3 target=subsection reason=relative-heading-rank; base=h2

\subsection{第十四条 教规}

关于失足的候洗者，神圣而伟大的主教会议规定：他们仅经过三年作为听道者之后，应与候洗者一同祈祷。\par

% structure: p0035 source=h3 target=subsection reason=relative-heading-rank; base=h2

\subsection{第十五条 教规}

由于发生重大扰乱与分歧，兹规定：在某些地方违反教规的惯例必须完全废止；即主教、司铎或执事不得从一个城市转到另一个城市。若有任何人，在神圣而伟大的主教会议此规定之后，试图行此类事或继续此类做法，他的行为完全无效，并应恢复他原被祝圣为主教或司铎的教会。\par

% structure: p0037 source=h3 target=subsection reason=relative-heading-rank; base=h2

\subsection{第十六条 教规}

凡司铎、执事或任何其他列入圣职名册者，若不敬畏天主，不遵守教会教规，竟轻率地离开自己的教会，绝不应被另一教会接纳；而应施加各种约束使其回归本堂区；若他们不愿去，则必须被绝罚。若有人胆敢私自带走属于另一教会的人，并在未经其本主教同意的情况下授圣职于自己的教会（虽然此人已被列入圣职名册，但已脱离本教会），则该授职无效。\par

% structure: p0039 source=h3 target=subsection reason=relative-heading-rank; base=h2

\subsection{第十七条 教规}

由于许多列入圣职名册者，贪图财物与利益，忘记了圣经的教训：\BibleQuote{他没有将钱放贷取利，}在借贷时索取百分之一的月息，神圣而伟大的主教会议认为公正的是：若在此规定之后，有人被发现收取利息，无论是通过秘密交易或其它方式，例如要求偿还本金外加一半，或为不义之财使用任何其他手段，则应将他从圣职中撤除，并将其名字从名册中删去。\par

% structure: p0041 source=h3 target=subsection reason=relative-heading-rank; base=h2

\subsection{第十八条 教规}

神圣大公会议获悉，在某些地区和城市，执事将圣体分送给司铎，然而，无论是典章还是习俗，均不允许无权奉献者将基督的圣体给予有权奉献者。此外，已知某些执事甚至在主教之前触摸圣体。凡此种种，均应彻底杜绝；执事应各守本位，自知是主教的仆役，且在司铎之下。他们应按其品级，在司铎之后领受圣体；或由主教，或由司铎分送给执事。再者，执事不应列坐于司铎之中，这违反典章和秩序。若在此法令之后，有人拒不服从，应罢免其执事之职。\par

% structure: p0043 source=h3 target=subsection reason=relative-heading-rank; base=h2

\subsection{第十九条法令}

关于投奔天主教会寻求庇护的保利安派，已颁布决议：他们必须重新受洗；若其中过去曾列于其司铎行列的人被发现无可指责、无可非议，则他们应由天主教会的主教重新受洗并祝圣；但若审查发现他们不适宜，则应予以罢免。同样，对于他们的女执事，以及一般对于曾列于其司铎行列的人，应遵守相同格式。我们所说的女执事，是指那些已穿上会服的人，但由于她们没有按手，只应算作平信徒。\par

% structure: p0045 source=h3 target=subsection reason=relative-heading-rank; base=h2

\subsection{第二十条法令}

\Emph{鉴於}有人在主日和五旬节期间跪下祈祷，因此为使一切在各地（各堂区）统一遵守，圣公会议认为，应站立向天主祈祷。\par

% structure: p0047 source=h2 target=section reason=relative-heading-rank; base=h2

\section{公会议信函}

致\Place{亚历山大}教会，赖天主的恩宠，神圣而伟大；并致我们在\Place{埃及}、\Place{五城地区}、\Place{利比亚}及天下万国的心爱的弟兄、正统的圣职人员及平信徒：在\Place{尼西亚}集会的神圣大公会议众主教，在主内祝你们安康。\par

\Emph{鉴於}在\Place{尼西亚}集会的伟大神圣的公会议，藉著基督的恩宠和我们最虔诚的君主\Person{君士坦丁}（他从各省各城召集我们），审议了有关教会信仰的事宜，我们以为有必要将若干事项以书面传达给你们，好使你们得以知道所讨论和调查的内容，以及所颁布和确认的决议。\par

首先，在我们最虔诚的君主\Person{君士坦丁}面前，调查了有关\Person{亚流}及其追随者的不敬和过犯；一致决定诅咒他和他的不敬观点，以及他所沉迷的亵渎之言和臆测，他亵渎天主圣子，说他从无有而来，在他被生之前他不存在，有过一段他不存在的时间，天主圣子凭其自由意志能犯恶也能行善，又说他是受造物。圣公会议诅咒了这一切，甚至无法容忍听他的不敬教义、狂妄和亵渎之言。至于对他的指控和后果，你们或已听闻或将了解详情，免得我们似乎是在压迫一个实际上已经因其罪过而受到适当报应的人。他的不敬甚至如此猖獗，以致他竟毁了\Place{马莫里卡}的\Person{托欧纳斯}和\Place{托勒迈斯}的\Person{塞孔德斯}，因为他们也接受了与其他相同的判决。\par

然而当天主的恩宠将\Place{埃及}从那种异端和亵渎中拯救出来，并从那些胆敢在向来和睦的百姓中制造骚乱和分裂的人中拯救出来，仍剩下关于\Person{梅里提乌}及其被祝圣者之傲慢的问题。关于这项工作的这一部分，我们现在，心爱的弟兄们，将告知你们公会议的决议。公会议于是决定温和对待\Person{梅里提乌}（因为按严格正义他不应得到宽恕），决议他应留在自己的城市，但无祝圣、管理或任命的权力，并不得为此目的出现在乡下或任何其他城市，但只享有其品级的空头衔。但那些由他安置的人，经更神圣的按手礼确认后，应按以下条件获准共融：他们应保留其品级并有权行圣职，但应完全低于所有在任何教会或堂区登记并已被我们最可敬的同僚\Person{亚历山大}任命的人。因此，未经我们最神圣的同僚\Person{亚历山大}手下的天主教会和宗徒教会的主教同意，这些人无权任命合自己心意的人，或提名，或做任何事。而那些赖天主的恩宠和你们的祈祷，未陷入分裂，反而在天主教会和宗徒教会中无瑕疵的人，有权在司铎中任命和提名合适的人，并有权按教会的法律和规定做一切事。但如果现在教会中的某些司铎去世，那么近期被接纳的人应继任亡者的职务，但必须他们显得合适，民众选举他们，\Place{亚历山大}主教同意选举并批准。这项让步已给予其余一切人；但鉴于\Person{梅里提乌}从一开始的犯乱行为及其性格的鲁莽急躁，并未对他本人做出同样的决议；反而，由于他是一个可能再次犯同样混乱的人，不得给予他任何权力或特权。\par

这些是对\Place{埃及}和至圣的\Place{亚历山大}教会特别关心的细节；但如果在我们最尊贵的主、我们的同僚与弟兄\Person{亚历山大}在场时，还有其他通过教规或其他法令颁布的事项，他本人将更详细地传达给你们，因为他在已完成的工作中既是引导者也是同工。\par

我们进一步向你们宣告关于神圣复活节协议的喜讯，即这一细节也已通过你们的祈祷得到妥善解决；以致我们所有在东方、从前遵循犹太人惯例的弟兄，从此要与罗马人、你们以及所有从起初就遵守复活节的人，在同一时间庆祝上述最神圣的复活节庆典。\par

因此，因着这些有益的成果，以及我们共同的和平与和谐，并因着一切异端的剪除，你们要以更大的尊荣和更深的爱德，接纳我们的同僚、你们的主教\Person{亚历山大}，他因与我们同在而使我们喜乐，且在如此高龄不辞辛劳，为使你们与我们众人之间建立和平。也请为我们众人祈祷，使所认为适宜的事得以坚立；因为我们相信，这些事是为悦纳全能天主及其独生子、我们的主耶稣基督，以及圣神而行的，愿光荣归于祂，直至永远。阿们。\par

\SourceRecord{Translated by Henry Percival.；Nicene and Post-Nicene Fathers, Second Series；Vol. 14.；Edited by Philip Schaff and Henry Wace.；Buffalo, NY: Christian Literature Publishing Co.,；1900.；<http://www.newadvent.org/fathers/3801.htm>.}
